From the attached table it can be read that northern galactic pole has the
following coordinates:
\[
  \alpha_{\mathrm{G}} = 12^{\mathrm{h}}51^{\mathrm{m}};\quad
  \delta_{\mathrm{G}} = +27^\circ 08'.
\]
It means that at the sidereal time
$\theta = 0^{\mathrm{h}}51^{\mathrm{m}}$ this pole is in lower culmination at
the northern side and below the horizon. \hfill[2]

The angle between galactic equatorial plane and the horizon is equal to
$90^\circ - \delta_{\mathrm{G}} + 90^\circ - \varphi = 103^\circ 18'$ and this
angle is an analogue of $\varepsilon$ in the equation of the ecliptic.
Analogue of R.A.\ in the equation of the ecliptic is $90^\circ - A$ because
$A$ is clockwise while R.A.\ is counterclockwise and at the intersection of
the horizon plane and galactic equatorial azimuth is equal to $90^\circ$.
\hfill[2]

So we get
$\operatorname{tg}(h) = \sin(90^\circ - A)\,
  \operatorname{tg}(180^\circ - \varphi - \delta_{\mathrm{G}})$. So,
\[
  h = \operatorname{arc\,tg}\left(\cos A \cdot
        \operatorname{tg}(-\varphi - \delta_{\mathrm{G}})\right)
    = \operatorname{arc\,tg}\left(-\cos A \cdot
        \operatorname{tg} 76^\circ 42'\right).
\]
\hfill[4]

Checking:
\begin{align*}
  A &= 0, & h &= -76^\circ 42', \\
  A &= 90, & h &= 0^\circ.
\end{align*}
\hfill[2]
